\section{What the pictures and measurements mean}
\label{sec:objects}

A trajectory is the motion generated by Eq.~\eqref{eq:rossler}. Rather than
inspect its whole path, we can record where it crosses a chosen surface.
That surface is a \emph{Poincar\'e section}, and the rule taking one crossing
to the next is its return map. A suitable section reduces the continuous
three-dimensional motion to a two-dimensional map.

Writing the flow as $\Phi^t$ and the section as $\Sigma$, the first-return map
is
\begin{equation}
  P:U\subset\Sigma\longrightarrow\Sigma,\qquad
  P(q)=\Phi^{\tau(q)}(q),
\end{equation}
where $\tau(q)>0$ is the next admissible crossing time. In the
Barrio-section calculations, $\Sigma$ is the plane $x=x_*$ through the
small equilibrium, with $\dot{x}>0$. The recovered historical section uses
$y=y_*$ and $x<x_*$; the crossing orientation is specified by each experiment.
Those choices are part of the measurement, not interchangeable conventions.

A scalar return plot shows one coordinate of successive crossings, such as
$y_{n+1}$ against $y_n$. If the points lie close to a curve, we can describe
its rising and falling pieces. One turning point gives two monotone branches;
two turning points give three. These are called \emph{unimodal} and
\emph{bimodal}, respectively. A clean scalar plot is useful evidence, but does
not by itself prove that the full return dynamics are one-dimensional.

For a smooth flow, a regular return between transverse sections is locally
invertible: nearby points can be followed back to their previous crossing
\citep[Secs.~6.4 and 12.2]{teschl2012ordinary}. Thus a turning point in a
scalar plot is not a singularity of the full Poincar\'e map. For an observable
$s$ and a declared invariant set $A$, the plotted relation is
\[
  \mathcal R_s=\{(s(q),s(Pq)):q\in A\}.
\]
It need not be the graph of a function. A genuine scalar quotient would need
a specified projection $\pi:A\to I=\pi(A)$ and map $F:I\to I$ satisfying
\begin{equation}
  \pi\circ P|_A=F\circ\pi.
  \label{eq:quotient-condition}
\end{equation}
Narrow conditional spread in a finite sample does not establish this
identity or an invariant strong-stable foliation. We therefore call the
measured change a \emph{finite-horizon scalar-return classifier transition}.
An invariant topology transition remains a separate explanatory target.

\begin{table}[tbp]
\centering
\small
\renewcommand{\arraystretch}{1.2}
\begin{tabularx}{\textwidth}{@{}P{0.23\textwidth}X@{}}
\toprule
Term & Meaning in this study \\
\midrule
Periodic orbit & A motion that repeats after its smallest return time $T$. \\
Attractor & An invariant set approached by a neighborhood of initial states. \\
Chaotic saddle & An invariant chaotic set that nearby trajectories may visit
for a long time before escaping. \\
Floquet multiplier & A factor describing how a small perturbation changes
after one circuit of a periodic orbit. \\
Period doubling (flip) & A local event with a transverse multiplier at $-1$;
a daughter orbit may require two circuits before repeating. \\
Homoclinic orbit & A trajectory that leaves an equilibrium and approaches
that same equilibrium again asymptotically. \\
Reinjection & The way a returning trajectory is routed into another part
of the return domain. \\
\bottomrule
\end{tabularx}
\caption{A short vocabulary for reading the figures. The precise numerical
tests are collected in Appendix~\ref{supp:methods}.}
\label{tab:vocabulary}
\end{table}

We distinguish the physical flow period $T$, historical period label
$p_{\rm hist}$, and counts $N_{\Sigma_J}$ and $N_{\Sigma_B}$ on the historical
and Barrio sections. A family labeled $p_{\rm hist}=6$ below has
$N_{\Sigma_B}=8$; neither integer means $T=6$ or $T=8$. A grazing can change
a section count without changing the underlying periodic orbit. The section
and counting convention must therefore accompany each reported period.

There is also a strict mathematical limit on \emph{superstability}. For a
regular periodic orbit, Liouville's formula \citep{teschl2012ordinary} gives
\begin{equation}
\label{eq:liouville}
  \det D\Phi^T(q)=
  \exp\!\left(\int_0^T [a+x(t)-c]\,dt\right)>0.
\end{equation}
The autonomous neutral multiplier is one, so neither transverse multiplier
can vanish. A zero derivative in a scalar model is therefore not an exact
zero Floquet multiplier of the full flow. This distinction corrects one of
our earlier search interpretations.
It does not invalidate superstability as a property of a scalar reduced map;
it invalidates our former identification of its zero derivative with a zero
transverse flow multiplier.

\begin{table}[tbp]
\centering
\small
\renewcommand{\arraystretch}{1.15}
\begin{tabularx}{\textwidth}{@{}P{0.25\textwidth}>{\raggedright\arraybackslash}X>{\raggedright\arraybackslash}X@{}}
\toprule
Parameter-space object & Defining event & Current scope \\
\midrule
Hopf locus & An equilibrium eigenpair crosses the imaginary axis. &
Analytic equation and numerical checks. \\
Period-doubling locus & A periodic orbit has a transverse multiplier $-1$. &
Corrected local flip curves and daughter checks. \\
Homoclinic locus & An orbit leaves and returns asymptotically to an equilibrium. &
Numerical candidates; no existence enclosure or established hub connection. \\
Scalar branch classifier & A declared finite return projection changes its
resolved monotone-piece count. & Sampled bracket, dependent on representation
and observation horizon. \\
Double-critical locus & Two specified scalar critical points belong to one
periodic itinerary. & Center search unresolved; not a zero flow multiplier. \\
\bottomrule
\end{tabularx}
\caption{Different parameter-space objects answer different questions. A fold in the computed
flip locus is not thereby a fold of a homoclinic locus or the historical
topology-transition curve.}
\label{tab:distinct-loci}
\end{table}
